\section{Discussion}

\subsection{Why the Pareto Curve Misses the Feasible Region}

Paper~10's discussion argued the H1 success and H2 cost are
``parameterised by $\tau$'' and form a trade-off curve. The empirical
curve sits entirely outside the joint feasible region: at every
$\tau$, the K=50 cost is at least $4$~pp deeper than Paper~10's
$-2$~pp tolerance. The mechanism is structural. The detector's
firing decision at a given window is a step function of $\tau$ via
the inequality $|\Delta_w| \geq \tau$; in the simulated regime the
calibration-target shifts $|\Delta_w|$ are either ``large''
(typically $> 0.1$) or ``small'' (typically $< 0.05$), with little
mass in between. The detector therefore behaves approximately
binarily across the grid, and the K=50 cost is approximately
constant across the $\tau$ values that all map to the same
``mostly-fire'' or ``mostly-don't-fire'' regime.

\subsection{What Richer Detectors Would Need to Do}

The negative feasibility result motivates the next paper class. A
detector that can reach the feasible region must decouple the K=200
hazard-avoidance from the K=50 over-firing penalty. Three operational
levers are available:

\textbf{(i) Capacity-aware threshold.} A detector that uses different
$\tau$ values at different capacities can in principle achieve
hazard-avoidance at K=200 with a small $\tau$ while permitting more
lag-1 use at K=50 with a large $\tau$. This requires the deployment
context to expose the operating capacity to the detector.

\textbf{(ii) Multi-window evidence accumulation.} CUSUM and Bayesian
change-point detectors accumulate evidence across windows and can
fire only when the cumulative shift exceeds a confidence threshold,
reducing false-positive firing rates at K=50 where individual
$|\Delta_w|$ values exceed $\tau$ but do not indicate a sustained
shift.

\textbf{(iii) Decoupled fire/strategy choice.} A detector that fires
into multiple strategies (e.g., \textsf{lag1}, \textsf{trail3},
\textsf{fixed}) rather than the binary lag-1/fixed choice could in
principle find a middle weight set that mitigates the K=50 cost
while preserving K=200 safety. Paper~8 falsified the simple trail3
substitute at single $\tau$; the multi-substitute variant is
distinct and untested.

These three directions are mutually compatible. The pre-registration
discipline that constrained Papers~7--11 requires that any such
richer detector be evaluated against locked hypotheses analogous to
those used here.

\subsection{The H1 Robustness Result}

The substantive H1 finding --- worst-case K=200 hazard is uniformly
$0.0$ across the entire $\tau$ grid --- is stronger than Paper~10's
single-$\tau$ H1 and deserves emphasis. The simple detector is
\textit{robustly} safe at K=200: regardless of threshold choice, no
adaptive procedure in this family produces a per-window loss at high
capacity. This is the operationally important positive that survives
the H3 rejection.

\subsection{The Pre-Registration Sign Errors}

The H2 and H4 sign reversals are an honest pre-registration mistake.
The protocol predicted that ``larger $\tau$ fires more'' (H4) and
that ``larger $\tau$ penalises K=50 more'' (H2), encoding an
intuition that does not survive a moment's mathematical reflection:
$|\Delta_w| \geq \tau$ becomes \textit{harder} to satisfy as $\tau$
grows. The decision rules were locked in this incorrect direction
and the data falsified them with $\rho = -0.97$ (H4) and $\rho =
+0.87$ (H2).

The pre-registration discipline produces this kind of mistake from
time to time: a locked-too-early sign assumption shows up as a
crisp data-direction rejection. The constructive interpretation is
that the data have falsified the protocol's stated direction and
the corrected direction is now the substantive scientific claim.
This is exactly the kind of negative result that the discipline is
designed to surface honestly.

\subsection{Operational Implication}

For an organisation deploying HygienePrio at fixed capacity in the
ranges studied, the calibration recipe is now operationally clear:
\begin{itemize}
\item At $K \leq 100$: deploy lag-1 online recalibration directly
(Paper~7's recovery ratio $\rho \approx 1.0$).
\item At $K \geq 200$: deploy Paper~4 fixed weights (Paper~6's
absolute floor stands, but no calibration strategy in the
Papers~7--11 sequence improves on it; adaptive switching avoids
harm but does not generate additional value).
\item Adaptive switching with $\tau \in [0.02, 0.10]$ adds
operational complexity without reaching the joint Paper~10
tolerances; richer detectors are required before adaptive procedures
should be preferred to the static rule.
\end{itemize}
